% !TEX TS-program = pdflatex
% Example: a Politecnico di Milano master's thesis in the "article"
% format, using the modular academicthesis class + academicthesis-polimi.
% Mirrors the official Scuola 3I template (Thesis_article_format).
% Finds the packages in ../core, ../languages, ../universities (any editor).
% academicthesis-repo-paths.tex -- lets showcases find the packages in this
% repository with ANY editor or engine (Texifier/Texpad, TeXShop, TeXstudio,
% plain pdflatex), not only with latexmk.
%
% \input it on the first line of a showcase, BEFORE \documentclass:
%
%   % academicthesis-repo-paths.tex -- lets showcases find the packages in this
% repository with ANY editor or engine (Texifier/Texpad, TeXShop, TeXstudio,
% plain pdflatex), not only with latexmk.
%
% \input it on the first line of a showcase, BEFORE \documentclass:
%
%   % academicthesis-repo-paths.tex -- lets showcases find the packages in this
% repository with ANY editor or engine (Texifier/Texpad, TeXShop, TeXstudio,
% plain pdflatex), not only with latexmk.
%
% \input it on the first line of a showcase, BEFORE \documentclass:
%
%   \input{academicthesis-repo-paths}
%   \documentclass[university=luiss]{academicthesis}
%
% It only extends LaTeX's own file search list (\input@path), which
% \documentclass, \RequirePackage, \input, \IfFileExists and
% \includegraphics all consult after TeX's normal search. The packages
% themselves stay path-free and are still loaded purely by name. Paths are
% relative to this showcases/ folder.
%
% \input@path does not search subfolders, so add one line per new
% university folder here (languages/ and core/ are flat).
\makeatletter
\providecommand*\input@path{}
\edef\input@path{%
  \input@path
  {../core/}%
  {../languages/}%
  {../universities/luiss/}%
  {../universities/polimi/}%
  {../universities/bocconi/}%
}
\makeatother
\endinput

%   \documentclass[university=luiss]{academicthesis}
%
% It only extends LaTeX's own file search list (\input@path), which
% \documentclass, \RequirePackage, \input, \IfFileExists and
% \includegraphics all consult after TeX's normal search. The packages
% themselves stay path-free and are still loaded purely by name. Paths are
% relative to this showcases/ folder.
%
% \input@path does not search subfolders, so add one line per new
% university folder here (languages/ and core/ are flat).
\makeatletter
\providecommand*\input@path{}
\edef\input@path{%
  \input@path
  {../core/}%
  {../languages/}%
  {../universities/luiss/}%
  {../universities/polimi/}%
  {../universities/bocconi/}%
}
\makeatother
\endinput

%   \documentclass[university=luiss]{academicthesis}
%
% It only extends LaTeX's own file search list (\input@path), which
% \documentclass, \RequirePackage, \input, \IfFileExists and
% \includegraphics all consult after TeX's normal search. The packages
% themselves stay path-free and are still loaded purely by name. Paths are
% relative to this showcases/ folder.
%
% \input@path does not search subfolders, so add one line per new
% university folder here (languages/ and core/ are flat).
\makeatletter
\providecommand*\input@path{}
\edef\input@path{%
  \input@path
  {../core/}%
  {../languages/}%
  {../universities/luiss/}%
  {../universities/polimi/}%
  {../universities/bocconi/}%
}
\makeatother
\endinput

\documentclass[university=polimi, master, format=article, language=english]{academicthesis}

% The class only loads what the PoliMi layout needs; anything else is
% your choice. These are the ones the official template example uses.
\usepackage{amssymb}
\usepackage{bm}
\usepackage{float}    % [H] placement
\usepackage{colortbl} % \rowcolor{bluepoli!40}
\usepackage[colorlinks=true, linkcolor=black, citecolor=black, urlcolor=black]{hyperref}

\thesissetup{
  title         = {Title},
  candidate     = {Name Surname},
  idnumber      = {Student ID},
  degree        = {Xxxxxxxxxxxx Engineering - Ingegneria Xxxxxxxxxxxx},
  supervisor    = {Prof. Name Surname},
  cosupervisors = {Name Surname, Name Surname},
  academicyear  = {20XX-20XX},
  % In the article format the English abstract and its keywords are
  % printed on the first page, so they are given here.
  abstract      = {Here goes the Abstract in English of your thesis (in
    article format) followed by a list of keywords. The Abstract is a
    concise summary of the content of the thesis (single page of text) and
    a guide to the most important contributions included in your thesis.
    It should be a self-contained text and should be clear to someone who
    hasn't (yet) read the whole manuscript.},
  keywords      = {here, the keywords, of your thesis},
}

%\addbibresource{example.bib}

\begin{document}

\frontpage

\section{Introduction}
This document is an example of the PoliMi article format built with the
modular \texttt{academicthesis} class. The first page carries the title,
the study programme, the author and student ID, the advisor, the
co-advisors, the academic year, the abstract and the keywords, all set
with \verb|\thesissetup|.

It is convenient to break the article format of your thesis into
sections and subsections. If necessary, subsubsections, paragraphs and
subparagraphs can be used.

\section{Equations}
Maxwell's equations read:
\begin{align}
  \nabla\cdot \bm{D} & = \rho, \label{eq:maxwell1} \\
  \nabla \times \bm{E} + \frac{\partial \bm{B}}{\partial t} & = \bm{0}, \\
  \nabla\cdot \bm{B} & = 0, \\
  \nabla \times \bm{H} - \frac{\partial \bm{D}}{\partial t} & = \bm{J}.
\end{align}

\section{Figures and tables}
\subsection{Figures}
\begin{figure}[H]
  \centering
  \includegraphics[width=0.3\textwidth]{logopolimi3i}
  \caption{Caption of the Figure.}
  \label{fig:logo}
\end{figure}

\subsection{Tables}
\begin{table}[H]
  \centering
  \begin{tabular}{|p{3em} c c c|}
    \hline
    \rowcolor{bluepoli!40}
    & \textbf{column 1} & \textbf{column 2} & \textbf{column 3} \\
    \hline \hline
    \textbf{row 1} & 1 & 2 & 3 \\
    \textbf{row 2} & $\alpha$ & $\beta$ & $\gamma$ \\
    \hline
  \end{tabular}
  \caption{Caption of the Table.}
  \label{table:example}
\end{table}

\section{Theorems and propositions}
\begin{theorem}
  Write here your theorem.
\end{theorem}
\textit{Proof.} If useful you can report here the proof.

\begin{proposition}
  Write here your proposition.
\end{proposition}

\section{Conclusions}
A final section containing the main conclusions of your research and
possible future developments of your work.

%\printbibliography

\appendix
\section{Appendix A}
If you need to include an appendix to support the research in your
thesis, you can place it at the end of the manuscript.

\clearpage

\begin{abstractthesis}[language=italian,
                       keywords={qui, le parole chiave, della tesi, in italiano}]
  Qui va l'Abstract in lingua italiana della tesi seguito dalla lista di
  parole chiave.
\end{abstractthesis}

\begin{acknowledgements}
  Here you might want to acknowledge someone.
\end{acknowledgements}

\end{document}
